\section{实用工具}

本章按比赛/训练的完整工作流组织：\textbf{编译}代码 $\to$ \textbf{计时}评估性能 $\to$ \textbf{生成测试数据}（先基础随机数据，再树/图等结构化数据，最后是 Python 方案）$\to$ \textbf{对拍}验证正确性。各节由前到后依次衔接，按需取用。

章末另附抗卡整数哈希，供比赛中使用哈希表时直接取用。

\subsection{Windows 自动编译器}

\lstinputlisting[language=bash]{codes/13-utilities-01.bat}

\textbf{使用方法}
\lstinputlisting[language=bash]{codes/13-utilities-02.bat}

\subsection{计时器}

利用 \code{std::chrono} 的 RAII 计时对象：构造时开始计时，作用域结束时自动输出耗时（毫秒）。配合上一节的编译器使用，可在本地快速评估程序实际运行时间；在 \code{main} 开头声明一个 \code{Timer} 即可测得整体耗时。

\lstinputlisting[language=C++]{codes/13-utilities-06.cpp}

\subsection{随机数生成器}

生成对拍所需的\textbf{基础随机数据}。C++11 起推荐使用 \code{mt19937}（梅森旋转引擎），比 \code{rand()} 周期长、分布均匀，且支持任意范围的均匀整数分布。

\lstinputlisting[language=C++]{codes/13-utilities-03.cpp}

\subsection{Python 随机数生成}

\textbf{用途：} 使用 \code{random} 标准库为对拍生成随机数据，或为随机化算法提供随机源。Python 写数据生成器比 C++ 更省事，适合规则复杂、规模不大的数据。

\begin{itemize}
  \item \code{random.seed(x)}：固定随机种子，使生成结果可复现；不传参则由系统熵源自动播种。
  \item \code{random} 模块是\textbf{确定性伪随机}（梅森旋转算法），输出序列可被预测；对抗 hack 的场景应改用 \code{random.SystemRandom()}。
  \item 批量生成大量随机数时，\code{random.choices} / \code{random.sample} 由 C 实现，比循环调用 \code{randint} 快数倍。
\end{itemize}

\lstinputlisting[language=Python]{codes/13-utilities-09.py}

\subsection{单文件对拍}

数据准备好后，用对拍验证正确性：把待测程序写成下面的 \code{Solution} 骨架，核心逻辑集中在 \code{solve()} 中，便于与暴力程序批量比对输出。

\lstinputlisting[language=C++]{codes/13-utilities-07.cpp}

\subsection{Python 对拍器}

批量对拍两个程序的输出，需要准备三个 C++ 源文件，脚本会先使用 GNU C++20、\code{-O2} 和 \code{-DONLINE\_JUDGE} 自动编译：
\begin{itemize}
    \item \code{gen.cpp}：随机数据生成器（可用本章前几节的生成器）
    \item \code{sol.cpp}：你的待测程序（由 \code{sol} 指定）
    \item \code{brute.cpp}：标准/暴力程序（由 \code{std} 指定）
\end{itemize}

在这三个源文件所在目录运行脚本，可修改 \code{sol}、\code{std}、\code{gen} 指定文件名。每次执行前打印命令；退出码非零时输出 Runtime Error、退出码、标准输出和标准错误并终止，单次编译或运行超过默认 $5$ 秒时捕获超时异常，打印 Time Limit Exceed 并终止。每轮将输入保存到 \code{in}，对输出调用 \code{rstrip()} 去除末尾空白后逐字符比较，内部空白仍参与比较；不一致时打印两份输出并停止，一致时打印绿色 AC。浮点容差或 Special Judge 题目需另写比较逻辑。

\lstinputlisting[language=Python]{codes/13-utilities-08.py}

\subsection{抗卡整数哈希 (Randomized Integer Hash)}

\textbf{用途：} 为 \code{unordered\_map} / \code{unordered\_set} 的整数键提供带随机盐的自定义哈希，降低预先构造的碰撞数据使哈希表退化、导致超时的风险。

\begin{itemize}
  \item \textbf{碰撞攻击：} 某些实现的默认整数哈希接近恒等映射，攻击者可按桶数构造大量落入同一桶的键；连续插入可能退化到 $O(n^2)$。
  \item \textbf{混合函数：} \code{SplitMix64} 先给输入加上 \code{seed}，再通过移位、异或和乘法混合整数各位；所有运算在 \code{uint64\_t} 上进行，溢出按模 $2^{64}$ 处理。
  \item \textbf{运行时随机盐：} 构造时混合 \code{steady\_clock}、\code{system\_clock} 的计数值与当前哈希器的对象地址，生成成员变量 \code{seed}，使攻击者难以提前确定映射。同一哈希器的种子在使用期间保持不变，默认拷贝会保留种子；独立构造的哈希器可以有不同种子。不能在容器使用期间修改种子，也不能每次调用时重新取时钟或随机数。
  \item \textbf{键的范围：} 适用于不超过 64 位的整数，包括 \code{int}、\code{long long} 和 \code{uint64\_t}；负数转换为无符号整数后按模 $2^{64}$ 处理。容器通过键的相等比较区分元素，哈希碰撞不会导致答案错误。
  \item \textbf{容量设置：} 先设置 \code{max\_load\_factor(0.7f)}，再用 \code{reserve(n)} 按预计不同键数量预留空间。两者减少扩容和平均桶长，但不能代替随机化哈希；\code{reserve} 后桶数仍可能变化，不要长期保存跨扩容的迭代器。
\end{itemize}

\textbf{时间复杂度：} 单次哈希 $O(1)$；分布良好时，查询、插入和删除平均 $O(1)$，插入扩容成本按均摊计，存储 $n$ 个不同键占用 $O(n)$ 空间。随机化不提供最坏情况保证，单次容器操作仍可能为 $O(n)$。

\begin{quote}
  \textbf{注意：} 这是竞赛中防预构造碰撞的实用措施，不是密码学哈希；时钟盐也不是安全随机数。需要稳定的最坏 $O(\log n)$ 上界或有序遍历时，使用 \code{map} / \code{set}。示例读入 $n$ 个有符号 64 位整数，统计频次并输出不同整数的数量；无符号 64 位键需将容器键类型和输入变量一并改为 \code{uint64\_t}。
\end{quote}

\lstinputlisting[language=C++]{codes/13-utilities-10.cpp}
